\vpara{Knowledge Base (KB).}
A KB is comprised by ontology $\{(C \times R \times C)\}$ and  relational facts $\{(E \times R \times (E\cup C))\}$, where $R,C,$ and $E$ denote relation set, class set, and entity set respectively.
Notably, we consider literal as a special type of entity.
Specifically, an ontology triple $(c_d, r, c_r)$ consists of a relation $r \in R$, a domain class $c_d$ which denotes the class of the subject entities, and a range class $c_r$ which denotes the class of the object entities.
Each class has multiple entities, thus an ontology triplet can be instantiated as several relational facts. For example, both $(e_1,r,e_2)$ and $(e_3,r,e_4)$ correspond to $(c_d, r, c_r)$, where $e_1, e_3 \in c_d$ and $e_2, e_4 \in c_r$.
Figure~\ref{fig:KB} illustrates a KB subgraph.
% We consider literal as a special type of entity. As shown in  
% Figure~\ref{fig:KB}, 

\vpara{SP-based KBQA.}
Given a natural question $q$, KBQA models aim to find a set of entities denoted by $A 
\subseteq E$ from KB as the answers to $q$. Instead of directly predicting $A$, SP-based KBQA models translate $q$ to an executable logical expression denoted by $s$ such as SPARQL, lambda-DCS~\cite{liang2013learning}, query graph~\cite{lan2020query}, and s-expression~\cite{gu2021beyond}. 

We select s-expression as our used logical expression since it could provide a good trade-off on compactness, compositionality, and readability~\cite{gu2021beyond}.
The \textbf{logical skeleton} of an s-expression can be derived by removing all the relations, classes, and entities in the expression and only keeping function operators and parentheses. 
Specifically, we replace relations, classes, entities, literals with special tokens ``<rel>'', ``<class>'', ``<entity>'', ``<literal>'' respectively. 
Figure~\ref{fig:KB} shows an executable logical expression on the KB and its corresponding logical skeleton. We unitedly name the relations, classes, entities, and logical skeleton in an s-expression as the \textbf{fine-grained component}, while the complete s-expression is the \textbf{coarse-grained logical expression}.

\begin{figure}
\centering
\includegraphics[scale=0.58]{knowledgegraph.png}
\caption{Illustration of a KB subgraph and an executable logical expression, where the ovals denote the entities, the rectangles denote the classes, the solid lines denote the relations, and the dashed lines connect the entities and their classes.
The upper part of the subgraph illustrates examples of ontology triplets, while the bottom illustrates relational facts.
}     
\label{fig:KB}
\end{figure}